\chapter{2010年上海卷（秋文）}

\section{填空题}

\begin{problem}
已知集合 \(A = \{1, 3, m\}\)，\(B = \{3, 4\}\)，\(A \cup B = \{1, 2, 3, 4\}\)，则 \(m =\fillinblank{}\).
\end{problem}

\begin{problem}
不等式 \(\frac{2 - x}{x + 4} > 0\) 的解集是\fillinblank{}.
\end{problem}

\begin{problem}
行列式 \(\left| \begin{array}{cc} \cos \frac{\pi}{6} & \sin \frac{\pi}{6} \\ \sin \frac{\pi}{6} & \cos \frac{\pi}{6} \end{array} \right|\) 的值是\fillinblank{}.
\end{problem}

\begin{problem}
若复数 \(z = 1 - 2\i\)（\(\i\) 为虚数单位），则 \(z \cdot \overline{z} + z =\fillinblank{}\).
\end{problem}

\begin{problem}
将一个总体分为 \(A\)、\(B\)、\(C\) 三层，其个体数之比为 \(5:3:2\). 若用分层抽样方法抽取容量为100的样本，则应从中抽取\fillinblank{} 个个体.
\end{problem}

\begin{problem}
已知四棱锥 \(P - ABCD\) 的底面是边长为 6 的正方形，侧棱 \(PA \perp\) 底面 \(ABCD\)，且 \(PA = 8\)，则该四棱锥的体积是\fillinblank{}.
\end{problem}

\begin{problem}
圆 \(C: x^{2} + y^{2} - 2x - 4y + 4 = 0\) 的圆心到直线 \(l: 3x + 4y + 4 = 0\) 的距离 \(d =\fillinblank{}\).
\end{problem}

\begin{problem}
动点 \(P\) 到点 \(F(2,0)\) 的距离与它到直线 \(x + 2 = 0\) 的距离相等，则点 \(P\) 的轨迹方程为\fillinblank{}.
\end{problem}

\begin{problem}
函数 \( f(x) = \log_3(x + 3) \) 的反函数的图象与 \( y \) 轴的交点坐标是\fillinblank{}.
\end{problem}

\begin{problem}
从一副混合后的扑克牌（52 张）中随机抽取 2 张，则“抽出的 2 张均为红桃”的概率为\fillinblank{}。（结果用最简分数表示）
\end{problem}

\begin{problem}
2010 年上海世博会园区每天 9：00 开园，20：00 停止入园. 在下边的框图中，\(S\) 表示上海世博会官方网站在每个整点报道的入园总人数，\(a\) 表示整点报道前 1 个小时内入园人数. 则空白的执行框内应填入\fillinblank{}.
\bitmapfigure[max width=0.42\linewidth,max totalheight=4.8cm]{img/2010/shanghai_liberal/q11_fig1.png}
\end{problem}

\begin{problem}
在 \(n\) 行 \(n\) 列矩阵 \(\left( \begin{array}{ccccccc}1 & 2 & 3 & \cdots & n - 2 & n - 1 & n\\ 2 & 3 & 4 & \cdots & n - 1 & n & 1\\ 3 & 4 & 5 & \cdots & n & 1 & 2\\ \vdots & \ddots & \vdots & \ddots & \vdots & \cdots & \vdots\\ n & 1 & 2 & \cdots & n - 3 & n - 2 & n - 1 \end{array} \right)\) 中，记位于第 \(i\) 行第 \(j\) 列的数为 \(a_{ij}\)，\(i, j = 1, 2, \cdots, n\). 当 \(n = 9\) 时， \( a_{11} + a_{22} + a_{33} + \cdots + a_{99} = \)\fillinblank{}.
\end{problem}

\begin{problem}
在平面直角坐标系中，双曲线 \(\Gamma\) 的中心在原点，它的一个焦点坐标为 \((\sqrt{5},0)\)，\(\bs{e}_1 = (2,1)\)，\(\bs{e}_2 = (2,-1)\) 分别是两条渐近线的方向向量. 任取双曲线 \(\Gamma\) 上的点 \(P\)，若 \(\overrightarrow{OP} = a\bs{e}_1 + b\bs{e}_2\)（\(a,b \in \R\)），则 \(a,b\) 满足的一个等式是\fillinblank{}.
\end{problem}

\begin{problem}
将直线 \(l_1: x + y - 1 = 0, l_2: nx + y - n = 0, l_3: x + ny - n = 0\) （\(n \in \mathbf{N}^*\)，\(n \geqslant 2\)）围成的三角形面积记为 \(S_n\). 则 \(\displaystyle\lim_{n \to \infty} S_n = \)\fillinblank{}.
\end{problem}

\section{单选题}

\begin{problem}
满足线性约束条件 \(\begin{cases} 2x + y \leqslant 3, \\ x + 2y \leqslant 3, \\ x \geqslant 0, \\ y \geqslant 0 \end{cases}\) 的目标函数 \(z = x + y\) 的最大值是（\quad）

\choices
    {1}
    {\(\frac{3}{2}\)}
    {2}
    {3}
\end{problem}

\begin{problem}
“\(x = 2k\pi + \frac{\pi}{4}\ (k \in \Z)\)”是“\(\tan x = 1\)”成立的（\quad）

\choices
    {充分不必要条件}
    {必要不充分条件}
    {充要条件}
    {既不充分也不必要条件}
\end{problem}

\begin{problem}
若 \(x_0\) 是方程 \(\lg x + x = 2\) 的解，则 \(x_0\) 属于区间（\quad）

\choices
    {(0,1)}
    {(1, 1.25)}
    {(1.25, 1.75)}
    {(1.75, 2)}
\end{problem}

\begin{problem}
若 \(\triangle ABC\) 的三个内角满足 \(\sin A: \sin B: \sin C = 5:11:13\)，则 \(\triangle ABC\)（\quad）

\choices
    {一定是锐角三角形}
    {一定是直角三角形}
    {一定是钝角三角形}
    {可能是锐角三角形，也可能是钝角三角形}
\end{problem}

\section{解答题}

\begin{problem}
已知 \(0 < x < \frac{\pi}{2}\)，化简：\(\lg\left(\cos x \cdot \tan x + 1 - 2\sin^2\frac{x}{2}\right) + \lg\left[\sqrt{2}\cos\left(x - \frac{\pi}{4}\right)\right] - \lg(1 + \sin 2x)\).
\end{problem}

\begin{problem}
如图所示，为了制作一个圆柱形灯笼，先要制作 4 个全等的矩形骨架. 总计耗用 9.6 米铁丝，再用 \(S\) 平方米塑料片制成圆柱的侧面和下底面（不安装上底面）.

\begin{enumerate}
\item 当圆柱底面半径 \(r\) 取何值时，\(S\) 取得最大值？并求出该最大值（结果精确到 0.01 平方米）；
\item 若要制作一个如图放置的，底面半径为 0.3 米的灯笼，请作出用于灯笼的三视图（作图时，不需考虑骨架等因素）.
\end{enumerate}

\bitmapfigure[max width=0.38\linewidth,max totalheight=4.8cm]{img/2010/shanghai_liberal/q20_fig1.png}

\end{problem}

\begin{problem}
\begin{enumerate}
\item 已知数列 \(\{a_{n}\}\) 的前 \(n\) 项和为 \(S_{n}\)，且 \(S_{n} = n - 5a_{n} - 85, n \in \mathbf{N}^{*}\)，证明：\(\{a_{n} - 1\}\) 是等比数列；
\item 求数列 \(\{S_{n}\}\) 的通项公式，并求出使得 \(S_{n+1} > S_{n}\) 成立的最小正整数 \(n\).
\end{enumerate}
\end{problem}

\begin{problem}
\begin{enumerate}
\item 若实数 \(x\)、\(y\)、\(m\) 满足 \(|x - m| < |y - m|\)，则称 \(x\) 比 \(y\) 接近 \(m\). 若 \(x^{2} - 1\) 比 \(3\) 接近 \(0\)，求 \(x\) 的取值范围；
\item 对任意两个不相等的正数 \(a, b\)，证明：\(a^{2}b + ab^{2}\) 比 \(a^{3} + b^{3}\) 接近 \(2ab\sqrt{ab}\)；
\item 已知函数 \(f(x)\) 的定义域 \(D = \{x \mid x \neq k\pi, k \in \Z, x \in \R\}\). 任取 \(x \in D\)，\(f(x)\) 等于 \(1 + \sin x\) 和 \(1 - \sin x\) 中接近 \(0\) 的那个值. 写出函数 \(f(x)\) 的解析式，并指出它的奇偶性、最小正周期、最小值和单调性（结论不要求证明）.
\end{enumerate}
\end{problem}

\begin{problem}
\begin{enumerate}
\item 已知椭圆 \(\Gamma\) 的方程为 \(\frac{x^2}{a^2} + \frac{y^2}{b^2} = 1\)（\(a > b > 0\)），\(A(0, b)\)，\(B(0, -b)\) 和 \(Q(a, 0)\) 为 \(\Gamma\) 的三个顶点. 若点 \(M\) 满足 \(\overrightarrow{AM} = \frac{1}{2}\left(\overrightarrow{AQ} + \overrightarrow{AB}\right)\)，求点 \(M\) 的坐标；
\item 设直线 \(l_{1}: y = k_{1}x + p\) 交椭圆 \(\Gamma\) 于 \(C, D\) 两点，交直线 \(l_{2}: y = k_{2}x\) 于点 \(E\). 若 \(k_{1} \cdot k_{2} = -\frac{b^{2}}{a^{2}}\)，证明：\(E\) 为 \(CD\) 的中点；
\item 设点 \(P\) 在椭圆 \(\Gamma\) 内且不在 \(x\) 轴上，如何作过 \(PQ\) 中点 \(F\) 的直线 \(l\)，使得 \(l\) 与椭圆 \(\Gamma\) 两个交点 \(P_{1}, P_{2}\) 满足 \(\overrightarrow{PP_{1}} + \overrightarrow{PP_{2}} = \overrightarrow{PQ}\)？令 \(a = 10\)，\(b = 5\)，点 \(P\) 的坐标是 \((-8, -1)\). 若椭圆 \(\Gamma\) 上的点 \(P_{1}, P_{2}\) 满足 \(\overrightarrow{PP_{1}} + \overrightarrow{PP_{2}} = \overrightarrow{PQ}\)，求点 \(P_{1}, P_{2}\) 的坐标.
\end{enumerate}
\end{problem}
